\documentclass[11pt]{article}
\usepackage{amsmath,amssymb,amsthm}

\newtheorem{theorem}{Theorem}
\newtheorem{corollary}[theorem]{Corollary}

\title{Cycle-lap rigidity for synchronized negative shadows}
\author{}
\date{27 August 2026}

\begin{document}
\maketitle

Let
\[
  2^{a_j}x_{j+1}=3x_j+1,
  \qquad
  2^{a_j}e_{j+1}=3e_j-1,
  \qquad a_j\geq1,
  \tag{1}
\]
where the $x_j$ are positive odd integers and the $e_j$ are real numbers
greater than $1$.  Put
\[
  m_j:=\lceil e_j\rceil,
  \qquad
  \theta_j:=m_j-e_j\in[0,1).
  \tag{2}
\]
Suppose that, on an interval of indices, there are a fixed
$c\in\{1,2,3\}$ and positive odd integers $y_j$ such that
\[
  m_j=cy_j,
  \qquad
  2^{a_j}y_{j+1}=3y_j-1.
  \tag{3}
\]
Thus $(y_j)$ is an accelerated positive $3y-1$ shadow synchronized with
the positive $3x+1$ path.

\begin{theorem}[synchronized cycle-lap rigidity]
Assume that (3) holds from $p$ through $p+KL$, where $K,L\geq1$, and that
\[
  y_{p+L}=y_p.
  \tag{4}
\]
Put
\[
  M:=y_p,
  \qquad
  B:=\sum_{i=0}^{L-1}a_{p+i},
  \qquad
  \rho:=\frac{3^L}{2^B}.
  \tag{5}
\]
Then the shadow and exponent blocks of length $L$ repeat throughout the
displayed interval, and
\[
\boxed{\begin{aligned}
  \rho&>1,\\
  v_2(x_p+M)&=KB+v_2(x_{p+KL}+M)\geq KB+1,\\
  x_{p+KL}+M&=\rho^K(x_p+M).
\end{aligned}}
\tag{6}
\]
Moreover,
\[
\boxed{
  \log\rho
  =\sum_{i=0}^{L-1}
    -\log\left(1-\frac1{3y_{p+i}}\right),
  \qquad
  \sum_{i=0}^{KL-1}\frac1{m_{p+i}}
  <\frac3c\log\frac{x_{p+KL}+M}{x_p+M}.}
\tag{7}
\]

The rounding error at the first return is exponentially pinched.  More
precisely,
\[
\boxed{
\begin{aligned}
  c=1&\quad\Longrightarrow\quad
    0\leq\theta_p<\rho^{-K},\\
  c=2,\ M=1&\quad\Longrightarrow\quad
    \rho=\frac32,
    \qquad 0<1-\theta_p\leq\left(\frac23\right)^K,\\
  (c-1)M>1&\quad\Longrightarrow\quad
    \rho^K<\frac{(c-1)M}{(c-1)M-1}.
\end{aligned}}
\tag{8}
\]

If $p$ is chosen at a maximum state $M$ of the shadow cycle, then
\[
  \rho\geq1+\frac1{2M}.
  \tag{9}
\]
Consequently,
\[
\boxed{
  c=3\Longrightarrow K\leq1,
  \qquad
  c=2,\ M\geq3\Longrightarrow K\leq2.}
\tag{10}
\]
Thus arbitrarily many complete laps are possible only for a color-$1$
shadow cycle or for the color-$2$ fixed shadow $y=1$.
\end{theorem}

\begin{proof}
Because $y_{j+1}$ is odd, the second identity in (3) determines both
\[
  a_j=v_2(3y_j-1)
  \qquad\text{and}\qquad
  y_{j+1}=\frac{3y_j-1}{2^{v_2(3y_j-1)}}.
  \tag{11}
\]
The shadow map is deterministic.  Equation (4) therefore makes both the
shadow states and the exponents $L$-periodic for as long as (3) holds.

Let $C$ be the positive affine carry obtained by iterating one period of
the negative recurrence.  We have
\[
  2^BM=3^LM-C,
  \qquad C>0.
  \tag{12}
\]
Hence $(3^L-2^B)M=C$, which proves $\rho>1$.  Adding the two synchronized
recurrences gives
\[
  2^{a_j}(x_{j+1}+y_{j+1})=3(x_j+y_j).
  \tag{13}
\]
After $K$ periods,
\[
  2^{KB}(x_{p+KL}+M)=3^{KL}(x_p+M).
  \tag{14}
\]
The two powers in (14) are coprime.  Since $x_{p+KL}$ and $M$ are odd,
taking $2$-adic valuations gives
\[
  v_2(x_p+M)=KB+v_2(x_{p+KL}+M)\geq KB+1.
\]
Division by $2^{KB}$ gives the last identity in (6).

For one shadow edge we have
\[
  \frac{y_{j+1}}{y_j}
  =\frac3{2^{a_j}}
   \left(1-\frac1{3y_j}\right).
  \tag{15}
\]
The product over one cycle equals $1$.  Taking logarithms proves the first
identity in (7).  Since $-\log(1-t)>t$ for $0<t<1$, it follows that
\[
  \sum_{i=0}^{L-1}\frac1{y_{p+i}}<3\log\rho.
  \tag{16}
\]
The same $L$ states occur in every lap and $m_j=cy_j$.  Thus
\[
  \sum_{i=0}^{KL-1}\frac1{m_{p+i}}
  =\frac Kc\sum_{i=0}^{L-1}\frac1{y_{p+i}}
  <\frac{3K}{c}\log\rho.
  \tag{17}
\]
Equation (6) turns the last member into the logarithmic ratio in (7).

Put $\Delta_j:=e_j-y_j$.  Subtracting the two negative recurrences gives
\[
  2^{a_j}\Delta_{j+1}=3\Delta_j,
  \qquad
  \Delta_{p+KL}=\rho^K\Delta_p.
  \tag{18}
\]
At a return to the state $M$, equations (2) and (3) give
\[
  (c-1)M-1<\Delta_j=(c-1)M-\theta_j\leq(c-1)M.
  \tag{19}
\]
If $c=1$, equation (18) says
$\theta_{p+KL}=\rho^K\theta_p<1$, which proves the first line of (8).
If $c=2$ and $M=1$, then the only positive shadow cycle is the fixed edge
$1\mapsto1$, with exponent $1$ and multiplier $3/2$.  Now (18)--(19) give
\[
  0<1-\theta_p
  =\rho^{-K}(1-\theta_{p+KL})
  \leq\rho^{-K}.
  \tag{20}
\]
Finally, if $(c-1)M>1$, then (18)--(19) give
\[
  \rho^K\bigl((c-1)M-1\bigr)<(c-1)M,
  \tag{21}
\]
which is the last line of (8).

We now rotate the cycle so that its displayed initial state $M$ is maximal.
If $a$ is the exponent on the edge returning to $M$, then
\[
  2^aM=3y-1\leq3M-1,
  \tag{22}
\]
where $y\leq M$ is the preceding cycle state.  Since $a\geq1$, equation
(22) forces $a=1$.  Writing the carry in (12) as
\[
  C=\sum_{i=0}^{L-1}3^{L-1-i}2^{B_i},
  \qquad
  B_i:=\sum_{h=0}^{i-1}a_{p+h},
  \tag{23}
\]
we see that its last normalized term is
$2^{B_{L-1}-B}=2^{-a}=1/2$.  Therefore
\[
  (\rho-1)M=\frac C{2^B}\geq\frac12,
  \tag{24}
\]
which proves (9).

Bernoulli's inequality and (9) give
\[
  \rho^K\geq1+K(\rho-1)\geq1+\frac K{2M}.
  \tag{25}
\]
For $c=3$, equations (8) and (25) imply
\[
  K<\frac{2M}{2M-1}<2,
\]
so $K\leq1$.  For $c=2$ and $M\geq3$, they imply
\[
  K<\frac{2M}{M-1}\leq3,
\]
so $K\leq2$.  This proves (10).
\end{proof}

\begin{corollary}[the two unbounded-lap mechanisms]
Consider a maximal constant nonpositive-defect run from the synchronized
ceiling construction, and write $c:=1-d\in\{1,2,3\}$.  Use its positive
odd shadow
\[
\begin{array}{c|c}
  c&\text{shadow states}\ \\ \hline
  1&y_j=m_j,\\
  2&y_j=m_j/2\quad\text{apart from the last ceiling state},\\
  3&y_j=m_j/3\quad\text{apart from the first ceiling state}.
\end{array}
\tag{26}
\]
Let $H:=10.789$ be the certified bound for the reciprocal mass of every
finite simple positive accelerated $3y-1$ path.  Then the following hold.

\begin{enumerate}
  \item A color-$3$ shadow traverses any one cycle at most once.  Its whole
  ceiling run has reciprocal mass less than
  \[
    H+\frac13.
    \tag{27}
  \]

  \item In a color-$2$ shadow, every cycle other than the fixed point $1$
  is traversed at most twice.  After deleting the terminal block on which
  $m_j=2$, the reciprocal mass of the ceiling run is less than
  \[
    2H+\frac12.
    \tag{28}
  \]
  If that terminal block contains $K$ edges and begins at $p$, then all of
  its exponents equal $1$ and
  \[
  \boxed{\begin{aligned}
    2^{K+1}&\mid x_p+1,\\
    x_{p+K}+1&=\left(\frac32\right)^K(x_p+1),\\
    0<1-\theta_p&\leq\left(\frac23\right)^K.
  \end{aligned}}
  \tag{29}
  \]

  \item In a color-$1$ run which contains two occurrences of a maximum
  state $M$ on its eventual shadow cycle, choose the first such occurrence
  $p$ and the last one $q=p+KL$.  The indices outside the complete laps
  $[p,q)$ form at most two simple shadow pieces.  Hence
  \[
  \boxed{\begin{aligned}
    \sum_{j\text{ in the run}}\frac1{m_j}
    &<2H+3\log\frac{x_q+M}{x_p+M},\\
    2^{KB+1}&\mid x_p+M,\\
    \theta_p&<\rho^{-K}.
  \end{aligned}}
  \tag{30}
  \]
  If no cycle maximum occurs twice, the run is covered by at most two simple
  shadow pieces and has mass less than $2H$.  If the complete shadow is
  simple, the sharper bound $H$ holds.
\end{enumerate}

Consequently, unbounded reciprocal mass inside individual nonpositive
defect runs can only be supplied by color-$1$ cycle laps or by terminal
color-$2$ towers with $m=2$ and exponent $1$.
\end{corollary}

\begin{proof}
A finite path under the deterministic accelerated $3y-1$ map is simple
until its first repeated state.  After that repetition it follows one cycle.
Every transient piece, every single traversal of the cycle with its repeated
endpoint removed, and every proper terminal piece is simple, so each has
reciprocal mass less than $H$.

For $c=3$, the theorem allows at most one complete cycle lap.  We split the
shadow into a simple initial piece, that lap, and a simple terminal piece.
Their total shadow mass is less than $3H$.  Division by $c=3$ and inclusion
of the discarded initial ceiling state, whose value is at least $3$, prove
(27).

For $c=2$, a nontrivial cycle has maximum at least $3$ and is traversed at
most twice.  Away from the fixed point $1$, the shadow is therefore covered
by a simple initial piece, at most two simple cycle laps, and a simple
terminal piece.  Its shadow mass is less than $4H$.  Division by $2$ and
inclusion of the discarded last ceiling state prove (28).  Once the shadow
reaches $1$, determinism gives the fixed edge $1\mapsto1$ with exponent
$1$.  Formula (29) is (6) and (8) for this fixed cycle.

For $c=1$, the shadow is the complete ceiling run.  If a cycle maximum
occurs at least twice, the portion before its first chosen occurrence and
the portion beginning at its last chosen occurrence are simple.  Their
combined mass is less than $2H$.  Apply (6)--(8) to the complete laps
between them to obtain (30).  If no cycle maximum occurs twice, splitting
at the first repeated state covers the path by at most two simple pieces.
The completely simple case is the definition of $H$.
\end{proof}

\begin{theorem}[fixed-seed plateau shell for the exceptional towers]
Assume now that (1) is the fixed synchronized companion of a normalized
divergent orbit.  Put
\[
  A_j:=\sum_{i<j}a_i,
  \qquad
  q_j:=\frac{3^j}{2^{A_j}},
  \tag{31}
\]
and let $\mathcal K$ be its stationary constant, so that
\[
  x_j+e_j=q_j\mathcal K.
  \tag{32}
\]
Define the centered rational anchor and its reduced ternary depth by
\[
  \Xi_j:=\frac{x_j+1}{q_j}
       =\frac{2^{A_j}(x_j+1)}{3^j},
  \qquad
  D_j:=\max\{0,j-v_3(x_j+1)\}.
  \tag{33}
\]
Then
\[
\boxed{
  \Xi_{j+1}-\Xi_j=\frac{2^{a_j}-2}{3q_j}\geq0,
  \qquad
  \mathcal K-\Xi_j=\frac{e_j-1}{q_j}>0.}
\tag{34}
\]
In particular, $\Xi_j$ is nondecreasing and changes exactly when $a_j>1$.

Let $[p,p+K]$ be an $m=2$ tower of $K\geq1$ edges, meaning that
\[
  m_{p+i}=2,
  \qquad
  a_{p+i}=1
  \qquad(0\leq i<K),
  \tag{35}
\]
and that $m_{p+K}=2$.  Put
\[
  b:=v_2(x_p+1),
  \qquad
  v:=v_3(x_p+1),
  \qquad
  h:=\frac{x_p+1}{2^b3^v}.
  \tag{36}
\]
Then $b\geq K+1$, $(h,6)=1$, and, for $0\leq i\leq K$,
\[
\boxed{
  x_{p+i}+1=2^{b-i}3^{v+i}h,
  \qquad
  \Xi_{p+i}=\Xi_p,
  \qquad
  D_{p+i}=D_p.}
\tag{37}
\]
The reduced denominator of $\Xi_p$ is exactly $3^{D_p}$.

Distinct maximal $m=2$ towers have distinct, strictly increasing anchors.
Consequently, for every integer $D\geq0$, the complete fixed orbit contains
at most
\[
\boxed{\left\lceil\mathcal K3^D\right\rceil-1}
\tag{38}
\]
maximal towers whose reduced ternary depth is at most $D$.

Consider two consecutive maximal towers.  Let their anchors be
$\xi_i<\xi_{i+1}$ and their depths be $D_i,D_{i+1}$.  Let $r_i$ be the
state index at which the first tower ends.  Then
\[
\boxed{
  3^{-\max\{D_i,D_{i+1}\}}
  \leq\xi_{i+1}-\xi_i
  <\mathcal K-\xi_i
  \leq\frac1{q_{r_i}}.}
\tag{39}
\]
Hence
\[
\boxed{
  \max\{D_i,D_{i+1}\}>\log_3 q_{r_i}.}
\tag{40}
\]

Finally, a tower beginning at $p$ satisfies the two fixed-seed length
bounds
\[
\boxed{
  K+1+v_3(x_p+1)\log_2 3
  <\log_2q_p+\log_2\mathcal K,}
\tag{41}
\]
and, with $n:=x_0$,
\[
\boxed{
  K+1+v_3(x_p+1)\log_2 3
  \leq p\log_2\frac32+\log_2(n+1).}
\tag{42}
\]
Thus an infinite family of exceptional towers must escape every bounded
reduced ternary shell, and every adjacent pair must expose a shell at least
as deep as the current coefficient surplus measured in base $3$.
\end{theorem}

\begin{proof}
Put $X_j:=x_j+1$.  The positive recurrence in (1) gives
\[
  2^{a_j}X_{j+1}=3X_j+2^{a_j}-2.
  \tag{43}
\]
After division by $2^{a_j}q_{j+1}=3q_j$, this becomes the first identity
in (34).  The second follows by subtracting $X_j=x_j+1$ from (32).

On the tower, equation (43) reduces to
\[
  X_{j+1}=\frac32X_j.
  \tag{44}
\]
Iteration proves the first identity in (37).  Since $X_{p+K}$ is even, its
$2$-adic valuation $b-K$ is at least $1$, so $b\geq K+1$.  The definition
of $h$ gives $(h,6)=1$.  Along the tower we also have
\[
  A_{p+i}=A_p+i,
  \qquad
  v_2(X_{p+i})=b-i,
  \qquad
  v_3(X_{p+i})=v+i.
  \tag{45}
\]
Substitution in (33) proves the other two identities in (37).  Cancellation
of the exact power $3^{\min\{p,v\}}$ in (33) shows that the remaining
denominator is precisely $3^{D_p}$.

Suppose that two maximal towers occur in this order.  If every intervening
exponent were $1$, then (1) would give
\[
  e_{j+1}-1=\frac32(e_j-1)
  \tag{46}
\]
throughout the intervening interval.  Once this increasing sequence leaves
the ceiling cell $1<e_j\leq2$, it cannot return.  Hence some intervening
exponent is greater than $1$.  Equation (34) proves that the second anchor
is strictly larger than the first.

If $D_j\leq D$, then the reduced anchor can be written as a positive integer
multiple of $3^{-D}$.  All anchors lie strictly between $0$ and $\mathcal K$
by (34).  Their strict increase therefore permits at most
$\lceil\mathcal K3^D\rceil-1$ such values.  This proves (38).

Two distinct rationals with reduced denominators $3^{D_i}$ and
$3^{D_{i+1}}$ differ by at least
$3^{-\max\{D_i,D_{i+1}\}}$.  At the endpoint $r_i$ of the first tower,
its anchor is still $\xi_i$, while $m_{r_i}=2$ gives
\[
  0<e_{r_i}-1\leq1.
  \tag{47}
\]
The second identity in (34) now proves every inequality in (39).  Taking
reciprocals of the two outer bounds gives (40).

The coprime divisors in (36) and the inequality $b\geq K+1$ give
\[
  2^{K+1}3^v\leq X_p.
  \tag{48}
\]
On the other hand, (34) gives $X_p<q_p\mathcal K$.  Taking binary
logarithms proves (41).

It remains to give a seed-only upper bound.  From (43), for every $a_j\geq1$
and $X_j\geq2$ we have
\[
\begin{aligned}
  X_{j+1}
  &=\frac{3X_j+2^{a_j}-2}{2^{a_j}}\\
  &\leq\frac32X_j,
\end{aligned}
\tag{49}
\]
because equality holds when $a_j=1$, while for $a_j\geq2$ the difference
after clearing denominators is
$(2^{a_j-1}-1)(3X_j-2)\geq0$.  Thus
\[
  X_p\leq\left(\frac32\right)^p(n+1).
  \tag{50}
\]
Combining (48) and (50), then taking binary logarithms, proves (42).
\end{proof}

\begin{theorem}[fixed-seed plateau shell for zero-defect laps]
Retain the normalized fixed-seed setup of the preceding theorem.  Put
\[
  N_j:=x_j+m_j,
  \qquad
  d_j:=2^{a_j}m_{j+1}-3m_j+1,
  \tag{51}
\]
and define the upper rational anchor and its reduced ternary depth by
\[
  \Omega_j:=\frac{N_j}{q_j}
  =\frac{2^{A_j}N_j}{3^j},
  \qquad
  E_j:=\max\{0,j-v_3(N_j)\}.
  \tag{52}
\]
Then
\[
\boxed{
  \Omega_{j+1}-\Omega_j=\frac{d_j}{3q_j},
  \qquad
  \Omega_j-\mathcal K=\frac{\theta_j}{q_j}\geq0.}
\tag{53}
\]
Every maximal zero-defect interval is finite.  On such an interval,
$\Omega_j$ and $E_j$ are constant, and the reduced denominator of its
anchor is exactly $3^{E_j}$.

For every integer $E\geq0$, only finitely many maximal zero-defect
intervals have reduced ternary depth at most $E$.

Let two maximal zero-defect intervals have distinct anchors
$\omega_i,\omega_{i+1}$ and depths $E_i,E_{i+1}$.  For each interval, let
\[
  Q_i:=\max\{q_j:j\text{ is a state index in that interval}\}.
  \tag{54}
\]
Then
\[
\boxed{
  3^{-\max\{E_i,E_{i+1}\}}
  \leq|\omega_{i+1}-\omega_i|
  <\frac1{Q_i}+\frac1{Q_{i+1}}
  \leq\frac2{\min\{Q_i,Q_{i+1}\}}.}
\tag{55}
\]
Consequently,
\[
\boxed{
  \max\{E_i,E_{i+1}\}
  >\log_3\frac{\min\{Q_i,Q_{i+1}\}}2.}
\tag{56}
\]

Suppose that one zero-defect interval contains $J\geq1$ complete laps of
a positive odd $3y-1$ cycle.  Choose a maximum cycle state $M$, let the
cycle have length $L$ and total exponent $B$, and let $p$ be the first
selected occurrence of $M$.  Then
\[
\boxed{
  JB+1+v_3(x_p+M)\log_2 3
  <\log_2q_p+\log_2(\mathcal K+1).}
\tag{57}
\]
Thus repeated zero-defect laps also escape every bounded ternary shell,
and their total exponent is paid by the coefficient surplus at their moving
start.
\end{theorem}

\begin{proof}
Adding the positive recurrence in (1) to the ceiling recurrence obtained
from (51) gives
\[
  2^{a_j}N_{j+1}=3N_j+d_j.
  \tag{58}
\]
Division by $2^{a_j}q_{j+1}=3q_j$ proves the first identity in (53).
The identities $N_j=x_j+m_j$ and $x_j+e_j=q_j\mathcal K$ prove the second.

When $d_j=0$, equation (58) shows that $\Omega_{j+1}=\Omega_j$ and
\[
  v_3(N_{j+1})=v_3(N_j)+1.
  \tag{59}
\]
Hence $E_{j+1}=E_j$.  Cancellation in (52), exactly as in the proof of
(37), shows that the reduced denominator is $3^{E_j}$.  A zero-defect tail
is impossible: iteration of (58) would force its unbounded accumulated
exponent to divide the fixed nonzero integer at the initial state.  Thus
every maximal zero-defect interval is finite.

We next prove finiteness in a bounded shell.  Since the normalized orbit
has $q_j>1$ for $j\geq1$ and $0\leq\theta_j<1$, equation (53) places every
upper anchor in the bounded interval
\[
  \mathcal K\leq\Omega_j<\mathcal K+1.
  \tag{60}
\]
There are only finitely many rationals in this interval whose reduced
denominator divides $3^E$.  Fix one such rational $R$.  If
$R>\mathcal K$, then (53) gives
\[
  q_j<\frac1{R-\mathcal K}
  \tag{61}
\]
at every occurrence of $R$.  Since $q_j\to\infty$, this happens only
finitely often.

It remains to consider $R=\mathcal K$.  If $\mathcal K$ is irrational,
this rational equality is impossible.  If $\mathcal K=u/v$ in lowest terms,
then an occurrence would make
\[
  q_j\mathcal K=\frac{u3^j}{v2^{A_j}}
  \tag{62}
\]
an integer.  The $2$-adic valuation of the right side tends to $-\infty$
because $A_j\to\infty$.  Thus (62) also has only finitely many solutions.
This proves bounded-shell finiteness.

Two distinct anchors with reduced denominators $3^{E_i}$ and
$3^{E_{i+1}}$ differ by at least the first member of (55).  On the other
hand, (53) and the definition of $Q_i$ give
\[
  0\leq\omega_i-\mathcal K<\frac1{Q_i},
  \qquad
  0\leq\omega_{i+1}-\mathcal K<\frac1{Q_{i+1}}.
  \tag{63}
\]
The triangle inequality proves the middle bound in (55), and the last bound
is immediate.  Rearrangement gives (56).

Finally, on a zero-defect cycle lap we have $m_j=y_j$.  At the selected
maximum state, $N_p=x_p+M$.  The cycle-lap theorem gives
\[
  2^{JB+1}\mid N_p.
  \tag{64}
\]
The coprime ternary divisor therefore gives
\[
  2^{JB+1}3^{v_3(N_p)}\leq N_p.
  \tag{65}
\]
Equations (52)--(53) and (60) give
\[
  N_p=q_p\Omega_p<q_p(\mathcal K+1).
  \tag{66}
\]
Taking binary logarithms in (65)--(66) proves (57).
\end{proof}

\begin{corollary}[bilateral exceptional-shell reduction]
Every repeated-mass interval left by the cycle-lap rigidity theorem has a
rational plateau anchor on one side of $\mathcal K$:
\[
\begin{array}{c|c|c}
\text{exceptional interval}&\text{anchor}&\text{side}\ \\ \hline
m=2\text{ exponent-one tower}&\Xi=(x+1)/q&\Xi<\mathcal K,\\
\text{zero-defect cycle laps}&\Omega=(x+m)/q&\Omega\geq\mathcal K.
\end{array}
\tag{67}
\]
The reduced denominator of either anchor is an exact power $3^D$ which is
constant on the interval.  Only finitely many exceptional intervals of
either kind lie in any bounded depth shell.

For an exceptional interval $I$, let $R_I$ be its anchor, $D_I$ its depth,
and $Q_I$ the maximum of $q_j$ on the interval.  If two such intervals have
distinct anchors, then
\[
\boxed{
  3^{-\max\{D_I,D_{I'}\}}
  \leq|R_I-R_{I'}|
  <\frac1{Q_I}+\frac1{Q_{I'}}.}
\tag{68}
\]
Hence every infinite family of exceptional repeated-mass intervals must
escape through unbounded ternary depth, quantitatively at the scale forced
by (40), (56), and (68).
\end{corollary}

\begin{proof}
The two anchor identities and the one-sided inequalities are (34) and (53).
The exact denominator and bounded-shell assertions are the preceding two
theorems.  Tower anchors are strictly increasing.  Moreover, on a tower
$q_j$ increases at every edge and
$\mathcal K-\Xi=(e_j-1)/q_j\leq1/q_j$, so two tower anchors satisfy (68)
directly.  For two distinct upper anchors, it is (55).  For one anchor of
each kind, $\Xi<\mathcal K\leq\Omega$, so
\[
  \Omega-\Xi
  =(\mathcal K-\Xi)+(\Omega-\mathcal K)
  <\frac1{Q_I}+\frac1{Q_{I'}}.
\]
Their power-of-$3$ denominators give the lower bound in (68).
\end{proof}

\begin{theorem}[universal constant-color run shell]
Retain the normalized fixed-seed setup, and put $c_j:=1-d_j$.  Let
$I=[s,t]$ be the state interval of a maximal constant nonpositive-defect
run, so that
\[
  c_j=c\in\{1,2,3\}
  \qquad(s\leq j<t).
  \tag{69}
\]
For $j\in I$, define
\[
\begin{aligned}
  Z_j(c)&:=cx_j+m_j,
  &\qquad
  \Lambda_I&:=\frac{Z_j(c)}{q_j},\\
  F_j(c)&:=\max\{0,j-v_3(Z_j(c))\},
  &
  F_I&:=F_j(c).
\end{aligned}
  \tag{70}
\]
These definitions are independent of the chosen $j\in I$.  More precisely,
\[
\boxed{
  2^{a_j}Z_{j+1}(c)=3Z_j(c),
  \qquad
  \frac{Z_{j+1}(c)}{q_{j+1}}=\frac{Z_j(c)}{q_j},
  \qquad
  F_{j+1}(c)=F_j(c).}
\tag{71}
\]
The reduced denominator of $\Lambda_I$ is exactly $3^{F_I}$.

The anchor lies on a prescribed side of $c\mathcal K$:
\[
\boxed{
  \Lambda_I-c\mathcal K
  =\frac{c\theta_j-(c-1)m_j}{q_j},}
\tag{72}
\]
and hence
\[
\begin{array}{c|c|c}
c&\text{side}&\text{distance bound valid for every }j\in I\\ \hline
1&\Lambda_I\geq\mathcal K&
  0\leq\Lambda_I-\mathcal K<1/q_j,\\
2&\Lambda_I<2\mathcal K&
  0<2\mathcal K-\Lambda_I\leq m_j/q_j,\\
3&\Lambda_I<3\mathcal K&
  0<3\mathcal K-\Lambda_I\leq2m_j/q_j.
\end{array}
\tag{73}
\]

For each fixed $c\in\{1,2,3\}$ and each integer $F\geq0$, only finitely
many color-$c$ runs have depth $F_I\leq F$.

For a run $I$, put
\[
  \varepsilon_I:=
  \begin{cases}
    \displaystyle\frac1{\max_{j\in I}q_j},&c=1,\\[3mm]
    \displaystyle\min_{j\in I}\frac{m_j}{q_j},&c=2,\\[3mm]
    \displaystyle2\min_{j\in I}\frac{m_j}{q_j},&c=3.
  \end{cases}
  \tag{74}
\]
If two color-$c$ runs $I,I'$ have distinct anchors, then
\[
\boxed{
  3^{-\max\{F_I,F_{I'}\}}
  \leq|\Lambda_I-\Lambda_{I'}|
  <\varepsilon_I+\varepsilon_{I'}.}
\tag{75}
\]
Moreover, $\varepsilon_I\to0$ along any sequence of runs whose initial
indices tend to infinity.  Thus every infinite family of constant
nonpositive-defect runs escapes through unbounded exact ternary depth.
\end{theorem}

\begin{proof}
The definition of $c_j$ rewrites the ceiling recurrence as
\[
  2^{a_j}m_{j+1}=3m_j-c_j.
  \tag{76}
\]
Multiplying the positive recurrence in (1) by $c_j=c$ and adding (76)
proves the first identity in (71).  Division by
$2^{a_j}q_{j+1}=3q_j$ proves the second.  The first identity also gives
\[
  v_3(Z_{j+1}(c))=v_3(Z_j(c))+1,
  \tag{77}
\]
which proves the third.  Cancellation in (70) shows that the exact reduced
denominator is $3^{F_I}$.  Taking $2$-adic valuations in the first identity
of (71) and iterating across the run gives
\[
  \sum_{j=s}^{t-1}a_j
  =v_2(Z_s(c))-v_2(Z_t(c))
  \leq v_2(Z_s(c)).
\]
Since every $a_j\geq1$, every constant-color run is finite.

Using $e_j=m_j-\theta_j$ and $x_j+e_j=q_j\mathcal K$, we obtain
\[
\begin{aligned}
  Z_j(c)
  &=c(x_j+e_j)+m_j-ce_j\\
  &=cq_j\mathcal K+c\theta_j-(c-1)m_j,
\end{aligned}
\]
which proves (72).  The bounds in (73) follow from
$0\leq\theta_j<1$ and $m_j\geq2$.

We have
\[
  \frac{m_j}{q_j}
  \leq\frac{e_j}{q_j}+\frac1{q_j}
  \longrightarrow0
  \tag{78}
\]
on the normalized ray.  Equations (72)--(73) therefore show that every
sequence of late color-$c$ anchors converges to $c\mathcal K$.  In
particular, all color-$c$ anchors lie in one bounded interval.

Fix a depth bound $F$.  Only finitely many rationals in that interval have
reduced denominator dividing $3^F$.  Fix one of them, say $R$.  If
$R\neq c\mathcal K$, convergence to $c\mathcal K$ shows that $R$ occurs
only finitely often.  If $R=c\mathcal K$, then either the equality is
impossible because $c\mathcal K$ is irrational, or $c\mathcal K=u/v$ is
rational in lowest terms.  In the latter case an occurrence would make
\[
  q_jc\mathcal K=\frac{u3^j}{v2^{A_j}}=Z_j(c)
  \tag{79}
\]
an integer.  Its $2$-adic valuation tends to $-\infty$ as
$A_j\to\infty$, so (79) also occurs only finitely often.  This proves
bounded-depth finiteness.

For two distinct anchors, their exact power-of-$3$ denominators give the
lower bound in (75).  The corresponding line of (73), minimized over each
run, gives the upper bound.  Finally, (78), together with $q_j\to\infty$
in the color-$1$ case, proves $\varepsilon_I\to0$ for late runs.  If an
infinite family had bounded depths, this would also contradict (75), in
agreement with the preceding finiteness argument.
\end{proof}

\begin{corollary}[all divergent nonpositive mass escapes every run shell]
On the normalized divergent ray, let $I(j)$ be the maximal constant-color
run containing a nonpositive-defect index $j$.  For every integer $F\geq0$,
\[
\boxed{
  \sum_{\substack{j\geq0:\ d_j\leq0\\F_{I(j)}>F}}
  \frac1{m_j}=\infty.}
\tag{80}
\]
Thus neither simple shadow fragments nor repeated shadow cycles can retain
divergent reciprocal mass in a bounded exact ternary shell.  The remaining
synchronized-shell problem is precisely to exclude divergent mass escaping
through $F_I\to\infty$ along one fixed ordinary orbit.
\end{corollary}

\begin{proof}
The total reciprocal mass on $d_j\leq0$ diverges.  By the theorem, only
finitely many constant-color runs have depth at most $F$.  Every one of
those runs is finite, so removing all of them subtracts a finite sum.  The
remaining sum is (80) and still diverges.
\end{proof}

\begin{theorem}[internal switches force maximal run depth]
Let two adjacent maximal nonpositive-defect runs meet at a state index
$s\geq1$, and let their distinct colors be
\[
  c,c'\in\{1,2,3\}.
  \tag{81}
\]
Let $\Lambda,\Lambda'$ and $F,F'$ be their anchors and depths from the
universal run-shell theorem, with the primed run entering at $s$.  Then
\[
\boxed{
  \Lambda'-\Lambda
  =\frac{(c'-c)x_s}{q_s}
  =\frac{(c'-c)2^{A_s}x_s}{3^s},}
\tag{82}
\]
and the last fraction has reduced ternary denominator exactly $3^s$.
More precisely,
\[
\boxed{F' =s,
\qquad F\leq s-1.}
\tag{83}
\]
Thus every late internal color switch places the newly entered run in the
deepest possible shell at that index.  If the new run ends at $t$, then
\[
\boxed{
  v_3(Z_j(c'))=j-s
  \qquad(s\leq j\leq t).}
\tag{83a}
\]
If it ends at another internal switch from $c'$ to $c''\neq c'$, put
\[
  r:=t-s,
  \qquad
  B:=\sum_{j=s}^{t-1}a_j.
\]
There is a positive integer $H$, not divisible by $3$, such that
\[
\boxed{
  Z_s(c')=2^BH,
  \qquad
  Z_t(c')=3^rH.}
\tag{83b}
\]
Moreover,
\[
  |c''-c'|=1\Longrightarrow H\text{ is odd},
  \qquad
  |c''-c'|=2\Longrightarrow H\text{ is even}.
\tag{83c}
\]
\end{theorem}

\begin{proof}
At the common state,
\[
  Z_s(c')-Z_s(c)=(c'-c)x_s.
  \tag{84}
\]
Division by $q_s$ proves (82).  The preceding positive recurrence gives
\[
  2^{a_{s-1}}x_s=3x_{s-1}+1\equiv1\pmod3,
  \tag{85}
\]
so $3\nmid x_s$.  Also
$c'-c\in\{-2,-1,1,2\}$, which is a unit modulo $3$.  The numerator in the
last member of (82) is therefore a unit modulo $3$, proving the denominator
claim.

Let the exiting run begin at $p<s$.  Iteration of (77) through this nonempty
run gives
\[
  v_3(Z_s(c))=v_3(Z_p(c))+s-p\geq1.
\]
Thus (84), together with the unit property of $(c'-c)x_s$, shows that
$Z_s(c')$ is a unit modulo $3$.  The entering depth is therefore $F'=s$.
The exiting depth is already defined at its initial state $p$, so
$F\leq p\leq s-1$.  This proves (83).  Finally, iteration of (77) from the
unit $Z_s(c')$ proves (83a).

Iteration of the first identity in (71) through the entering run gives
\[
  2^BZ_t(c')=3^rZ_s(c').
\]
By (83a), $Z_t(c')=3^rH$ for a positive integer $H$ with $3\nmid H$.
Cancellation gives $Z_s(c')=2^BH$, which proves (83b).  At the exit,
$Z_t(c'')$ is even because it begins a color-$c''$ edge, and
\[
  Z_t(c'')-Z_t(c')=(c''-c')x_t.
\]
The state $x_t$ is odd.  Thus $Z_t(c')$, and hence $H$, is odd when the
color difference is odd and even when it equals $2$ in absolute value.
This proves (83c).
\end{proof}

\begin{corollary}[transition-rich versus entrance-mass dichotomy]
On a normalized divergent ray, exactly one of the following two branches
holds.

\begin{enumerate}
  \item There are infinitely many internal switches between colors in
  $\{1,2,3\}$.  Then infinitely many adjacent run pairs satisfy the
  maximal-depth law (83).

  \item There are only finitely many such switches.  After deleting a
  finite prefix, every maximal nonpositive-defect block is one entrance
  run.  The divergent nonpositive reciprocal mass is then entrance-run
  mass.  Hence the run-mass valuation budget forces divergence of the sum
  of the preceding exponents or of the positive entrance-switch gains.
\end{enumerate}

In particular, no contradiction can follow from large run depth or large
power-of-$3$ anchor denominator alone: branch $1$ produces maximal depth
automatically.  A successful synchronized-shell proof must control
reciprocal mass relative to anchor width in branch $1$, or close the
preceding-exponent and switch-gain alternatives in branch $2$.
\end{corollary}

\begin{proof}
If the first branch holds, the theorem applies at every internal switch.
If it fails, only finitely many nonpositive blocks contain more than their
first constant-color run.  Removing those finite blocks leaves only entrance
runs.  Their total mass still diverges, and the stated alternatives are the
run-mass valuation budget applied and summed over these runs.  The final
assertion is the logical content of the two branches.
\end{proof}

\begin{theorem}[dyadic numerator law for a ternary run shell]
Retain the normalized fixed-seed setup.  Let $I=[s,t]$ be the state
interval of a constant nonpositive-defect run, with color
$c\in\{1,2,3\}$ and at least one edge.  Put
\[
  B_I:=A_t-A_s,
  \qquad
  U_I:=v_2(Z_s(c)),
  \qquad
  V_I:=v_2(Z_t(c)).
  \tag{86}
\]
For $s\leq j\leq t$, define
\[
  w_j:=\min\{j,v_3(Z_j(c))\},
  \qquad
  R_I:=\frac{2^{A_j}Z_j(c)}{3^{w_j}}.
  \tag{87}
\]
Then $R_I$ is a positive integer independent of $j$, and
\[
\boxed{
  \Lambda_I=\frac{R_I}{3^{F_I}},
  \qquad
  v_2(R_I)=A_j+v_2(Z_j(c)).}
\tag{88}
\]
In particular,
\[
\boxed{
  A_s+U_I=A_t+V_I,
  \qquad
  B_I=U_I-V_I.}
\tag{89}
\]

Write
\[
  R_I=2^{A_t+V_I}\eta_I,
  \qquad \eta_I\in\mathbb N\text{ odd}.
  \tag{90}
\]
The exponent mass of the run has the exact shell decomposition
\[
\boxed{
  B_I
  =F_I\log_2 3-A_s+\log_2\Lambda_I
   -V_I-\log_2\eta_I.}
\tag{91}
\]
If
\[
  \kappa_1=\kappa_3:=\frac13,
  \qquad
  \kappa_2:=\frac12,
  \tag{92}
\]
then
\[
\boxed{
  \sum_{j=s}^{t-1}\frac1{m_j}
  \leq\kappa_cB_I.}
\tag{93}
\]

Discard a finite prefix so that $q_s\geq1$, and put
$C:=3\mathcal K+1$.  Every remaining run satisfies
\[
\boxed{\begin{aligned}
  B_I
  &<\log_2(Cq_s)-(s-F_I)\log_2 3,\\
  \sum_{j=s}^{t-1}\frac1{m_j}
  &<\kappa_c\bigl(\log_2(Cq_s)
        -(s-F_I)\log_2 3\bigr),\\
  F_I
  &>s-\log_3(Cq_s)
   \geq s\log_3 2-\log_3 C.
\end{aligned}}
\tag{94}
\]
Thus linear escape in ternary depth is automatic on every late
nonpositive-defect run.  The quantity not controlled by depth alone is the
positive shell slack in the first line of (94).
\end{theorem}

\begin{proof}
Equation (77) and the identity $A_{j+1}=A_j+a_j$ give
\[
  w_{j+1}=w_j+1,
  \qquad
  \frac{2^{A_{j+1}}Z_{j+1}(c)}{3^{w_{j+1}}}
  =\frac{2^{A_j}Z_j(c)}{3^{w_j}}.
\]
This proves that $R_I$ is independent of $j$.  Since
$F_I=j-w_j$, the anchor identity in (88) follows.  Division by a power of
$3$ does not change a $2$-adic valuation, which proves the second identity
in (88).  Evaluating it at the two endpoints proves (89), and (90) follows.

We can compute the logarithm of $R_I$ in two ways:
\[
\begin{aligned}
  \log_2R_I
  &=F_I\log_2 3+\log_2\Lambda_I,\\
  \log_2R_I
  &=A_t+V_I+\log_2\eta_I
   =A_s+B_I+V_I+\log_2\eta_I.
\end{aligned}
\]
Their equality is exactly (91).

There are at most $B_I$ indices in the run because every exponent is at
least $1$.  On every displayed edge, parity in
\[
  2^{a_j}m_{j+1}=3m_j-c
\]
shows that $m_j$ is odd for $c\in\{1,3\}$ and even for $c=2$.
Thus $m_j\geq3$ in the first case and $m_j\geq2$ in the second.  This proves
(93).

The side bounds (73) and $q_s\geq1$ give
\[
  0<\Lambda_I<C.
\]
Since $V_I\geq0$ and $\eta_I\geq1$, equation (91) now gives
\[
  0<B_I
  <F_I\log_2 3-A_s+\log_2C.
\]
Using
\[
  A_s=s\log_2 3-\log_2q_s
\]
proves the first line of (94), and (93) proves the second.  Positivity of
$B_I$ gives $F_I>s-\log_3(Cq_s)$.  Finally,
$A_s\geq s$ implies $q_s\leq(3/2)^s$, and substitution proves the last
inequality in (94).
\end{proof}

\begin{corollary}[the entrance-exponent alternative is automatic]
Assume the finite-transition branch of the preceding dichotomy, and let
$\mathcal E$ be the family of all sufficiently late entrance runs.  Then
\[
  |\mathcal E|=\infty,
  \qquad
  \sum_{I\in\mathcal E}a_{s_I-1}=\infty.
  \tag{95}
\]
Consequently, divergence of the preceding-exponent sum in the run-mass
valuation budget is automatic and cannot distinguish a divergent ray.  A
necessary condition for the entrance-run mass to diverge is the sharper
shell-slack law
\[
\boxed{
  \sum_{I\in\mathcal E}
  \bigl(\log_2(Cq_{s_I})
       -(s_I-F_I)\log_2 3\bigr)=\infty.}
\tag{96}
\]
Thus the finite-transition branch has been reduced to controlling this
exact shell slack or to a strictly stronger compensated switch-gain bound.
\end{corollary}

\begin{proof}
Every entrance run is finite.  Since its reciprocal mass is finite, the
divergent total entrance-run mass requires infinitely many distinct entrance
runs.  Their preceding indices are distinct and every accelerated exponent
is at least $1$, which proves (95).  Summing the second inequality in (94)
over $\mathcal E$ proves (96), because
$\kappa_c\leq1/2$ for every $c\in\{1,2,3\}$.
\end{proof}

\begin{theorem}[rational barriers and logarithmic run transport]
Retain the normalized fixed-seed setup, and define
\[
  h_j:=\frac{x_j}{q_j},
  \qquad
  g_j:=\mathcal K-h_j=\frac{e_j}{q_j},
  \qquad
  u_j:=\frac{m_j}{q_j}.
  \tag{97}
\]
Then
\[
\boxed{\begin{aligned}
  h_{j+1}-h_j&=g_j-g_{j+1}=\frac1{3q_j},\\
  u_{j+1}-u_j&=-\frac{c_j}{3q_j},\\
  g_j&\leq u_j<g_j+\frac1{q_j}<2g_j.
\end{aligned}}
\tag{98}
\]
At a fixed state, introduce the three candidate rational barriers
\[
  \beta_j(c):=h_j+\frac{u_j}{c}
  =\frac{cx_j+m_j}{cq_j}
  \qquad(c\in\{1,2,3\}).
  \tag{99}
\]
They satisfy
\[
\boxed{
  \beta_j(3)<\beta_j(2)<\mathcal K\leq\beta_j(1),}
\tag{100}
\]
with exact gaps
\[
\begin{aligned}
  \beta_j(1)-\mathcal K&=\frac{\theta_j}{q_j},\\
  \mathcal K-\beta_j(2)&=\frac{2e_j-m_j}{2q_j},\\
  \mathcal K-\beta_j(3)&=\frac{3e_j-m_j}{3q_j}.
\end{aligned}
\tag{101}
\]

Let $I=[s,t]$ be a color-$c$ nonpositive-defect run.  Its barrier
\[
  \beta_I:=\frac{\Lambda_I}{c}
  \tag{102}
\]
is constant on the run, and
\[
\boxed{
  u_j=c(\beta_I-h_j),
  \qquad
  \frac{u_{j+1}}{u_j}=1-\frac{c}{3m_j}.}
\tag{103}
\]
Consequently,
\[
\boxed{
  \frac2c\log\frac{\beta_I-h_s}{\beta_I-h_t}
  \leq\sum_{j=s}^{t-1}\frac1{m_j}
  <\frac3c\log\frac{\beta_I-h_s}{\beta_I-h_t}.}
\tag{104}
\]

Two useful uniform consequences are
\[
\boxed{\begin{aligned}
  c=3&\quad\Longrightarrow\quad
  \sum_{j=s}^{t-1}\frac1{m_j}<\log2,\\
  c=2&\quad\Longrightarrow\quad
  \sum_{j=s}^{r-1}\frac1{m_j}<\frac32\log2,
\end{aligned}}
\tag{105}
\]
where in the second line $r$ is the first state of the run with $m_r=2$,
or $r=t$ if no such state occurs.  After the first state with $m=2$, every
remaining color-$2$ edge has exponent $1$ and remains at $m=2$.  Thus (105)
bounds the entire color-$2$ run outside its optional terminal tower.  If a
color-$2$ run ends at an internal switch, then the terminal tower is absent,
$r=t$, and consequently
\[
\boxed{
  \sum_{j=s}^{t-1}\frac1{m_j}<\frac32\log2.}
\tag{105a}
\]

Finally, suppose that an internal switch at state $s$ changes the color from
$c$ to $c'$.  The corresponding barrier jump obeys
\[
\boxed{
  \beta_s(c')-\beta_s(c)
  =u_s\left(\frac1{c'}-\frac1c\right),
  \qquad
  \frac{|\beta_s(c')-\beta_s(c)|}{m_s}
  =\left|\frac1{c'}-\frac1c\right|\frac1{q_s}.}
\tag{106}
\]
If $\mathcal S$ is the set of all internal nonpositive-color switches, then
\[
\boxed{
  \sum_{s\in\mathcal S}
  \frac{|\beta_s(c_s^+)-\beta_s(c_s^-)|}{m_s}
  \leq2(\mathcal K-h_0)<\infty.}
\tag{107}
\]
Hence the raw weighted barrier-jump charge is already summable.  Any
remaining transition-rich divergence must be carried by the logarithmic
barrier approaches in (104), not by the switch jumps in (106) themselves.
\end{theorem}

\begin{proof}
Division of the positive recurrence by
$2^{a_j}q_{j+1}=3q_j$ gives the first identity in (98).  The ceiling
recurrence gives the second.  Since $m_j=e_j+\theta_j$ and
$0\leq\theta_j<1<e_j$, the last line of (98) follows.

Equation (99) and $u_j=m_j/q_j$ give the ordering of the three candidate
barriers.  More precisely, $m_j=e_j+\theta_j$ gives (101).  The second and
third gaps are positive because $m_j<e_j+1<2e_j<3e_j$.

On a color-$c$ run, (71) says that $\Lambda_I$ is constant.  Equations
(99) and (102) therefore give the first identity in (103).  The second line
of (98) gives
\[
  \frac{u_{j+1}}{u_j}
  =1-\frac{c}{3q_ju_j}
  =1-\frac{c}{3m_j},
\]
which proves the second identity in (103).  Put
$z_j:=c/(3m_j)$.  Parity in the ceiling recurrence gives
$0<z_j\leq1/3$.  Hence
\[
  z_j< -\log(1-z_j)\leq\frac{z_j}{1-z_j}
  \leq\frac32z_j.
\]
Sum this inequality and use (103).  The result is (104).

Suppose first that $c=3$, and put
$\delta:=3\mathcal K-\Lambda_I>0$.  Along the run,
\[
  u_j=3g_j-\delta.
\tag{108}
\]
At the initial state, parity gives $m_s\geq3$.  Therefore
\[
  \frac{u_s}{g_s}=\frac{m_s}{e_s}
  <\frac{m_s}{m_s-1}\leq\frac32.
\]
Equation (108) now gives $u_s<\delta$.  At the terminal state,
$u_t\geq g_t$, so (108) gives $g_t\geq\delta/2$ and
$u_t\geq\delta/2$.  Thus $u_s/u_t<2$, and the first line of (105) follows
from the upper bound in (104).

Now let $c=2$.  If the prefix in (105) is nonempty, then $m_s\geq4$, so
\[
  \frac{u_s}{g_s}=\frac{m_s}{e_s}
  <\frac{m_s}{m_s-1}\leq\frac43.
\]
Put $\delta:=2\mathcal K-\Lambda_I>0$.  Since
$u_j=2g_j-\delta$, the last inequality gives $u_s<2\delta$.  At the
terminal state $r$ of the prefix, $u_r\geq g_r$ gives
$u_r\geq\delta$.  Therefore $u_s/u_r<2$, and the second line of (105)
follows from (104).  If $m_j=2$ on a color-$2$ edge, then
\[
  2^{a_j}m_{j+1}=4.
\]
Since $m_{j+1}\geq2$ and $a_j\geq1$, this forces
$a_j=1$ and $m_{j+1}=2$.  This proves the terminal-tower assertion.
If the run ends at an internal switch, the next color belongs to
$\{1,3\}$.  Reduction of (76) modulo $2$ gives
$m_t\equiv c_t\pmod2$, so $m_t$ is odd.  A color-$2$ tower would instead
force $m_t=2$.  Hence no such tower occurs and $r=t$, proving (105a).

At an internal switch the two barriers in (99) are evaluated at the same
point $(h_s,u_s)$, which proves (106).  The largest possible value of
$|1/c'-1/c|$ is $2/3$.  Summing the first identity in (98) gives
\[
  \sum_{j=0}^{\infty}\frac1{q_j}
  =3(\mathcal K-h_0).
\]
The second identity in (106), summed over the subset $\mathcal S$, now
proves (107).
\end{proof}

\begin{theorem}[fixed-ray carrier escape and collision law]
Let $I=[s,t]$ be an internally bracketed color-$c$ run on the normalized
divergent ray.  Put $r:=t-s$, and let $H$ be its carrier.  Then
\[
\boxed{
  \frac{H}{m_t}
  =\frac{1+c x_t/m_t}{3^r}.}
\tag{109}
\]
In particular,
\[
  H\leq m_t
  \quad\Longrightarrow\quad
  r\log3\geq\log\left(1+\frac{cx_t}{m_t}\right)
  >\log\frac{x_t}{m_t}.
\tag{110}
\]
Consequently, for every fixed $R$, all sufficiently late internally
bracketed runs of length at most $R$ satisfy $H>m_t$.

Now let $I'=[t,u]$ be the next run and suppose that it also ends at an
internal switch.  Write its color as $c'\neq c$, its exponent sum as
\[
  B':=\sum_{j=t}^{u-1}a_j,
\]
and its carrier as $H'$.  At the common state,
\[
\boxed{\begin{aligned}
  (c'-c)x_t&=2^{B'}H'-3^rH,\\
  (c'-c)m_t&=c'3^rH-c2^{B'}H'.
\end{aligned}}
\tag{111}
\]
If the adjacent carriers collide, $H'=H$, then
\[
\boxed{
  \frac{m_t}{x_t}
  =\frac{|c'3^r-c2^{B'}|}{|2^{B'}-3^r|}
  >\frac1{\max\{2^{B'},3^r\}},}
\tag{112}
\]
and hence
\[
  \max\{B'\log2,r\log3\}>\log\frac{x_t}{m_t}.
\tag{113}
\]
Thus, for every fixed $R$, only finitely many adjacent carrier collisions
can satisfy both $r\leq R$ and $B'\leq R$.  Every infinite collision family
on one divergent ray must escape through unbounded run length or unbounded
entering exponent sum.
\end{theorem}

\begin{proof}
The terminal identity in (83b) is
\[
  cx_t+m_t=3^rH.
\]
Division by $3^rm_t$ proves (109), and (110) follows immediately.  Equations
(32) and (78) give
\[
  \frac{m_t}{x_t}
  =\frac{m_t/q_t}{x_t/q_t}\longrightarrow0,
\]
because $x_t/q_t\to\mathcal K>0$.  This proves the bounded-length
consequence.

At the common state, the terminal identity for $I$ and the initial identity
for $I'$ are
\[
  cx_t+m_t=3^rH,
  \qquad
  c'x_t+m_t=2^{B'}H'.
\]
Subtraction gives the first line of (111).  Multiplying the first identity
by $c'$, the second by $c$, and subtracting gives the second line.

Suppose that $H'=H$.  The first line of (111) shows that
$2^{B'}-3^r\neq0$, while the second shows that
$c'3^r-c2^{B'}\neq0$.  After division, both factors $H$ and $c'-c$ cancel,
which proves the equality in (112).  Its numerator is a nonzero integer, so
its absolute value is at least $1$.  Its denominator is strictly smaller
than $\max\{2^{B'},3^r\}$.  This proves the strict inequality in (112) and,
after taking logarithms, (113).  The limit $m_t/x_t\to0$ proves the final
finiteness assertion.
\end{proof}

\begin{theorem}[finite strict-barrier carrier counterexamples]
Neither a terminal bound $H\geq m_t$ nor injectivity of the carriers of
adjacent internally bracketed runs follows from the synchronized
recurrences, the ceiling condition, and the strict-barrier inequalities.

More precisely, the finite data
\[
\begin{array}{c|rrrr}
j&0&1&2&3\\ \hline
x_j&91&137&103&155\\
m_j&48&71&53&78\\
\theta_j&26/27&17/18&17/24&1/16
\end{array}
\qquad
(a_0,a_1,a_2)=(1,2,1)
\]
give a synchronized real companion with color word $(2,1,3)$.  The middle
color-$1$ run is internally bracketed and has carrier
\[
  x_1+m_1=208=2^2\cdot52,
  \qquad
  x_2+m_2=156=3\cdot52.
\]
Thus $H=52<m_2=53$.

There is also a strict-barrier synchronized segment
\[
\begin{array}{c|rrrrrr}
j&0&1&2&3&4&5\\ \hline
x_j&27&41&31&47&71&107\\
m_j&29&43&32&47&70&104\\
\theta_j&1/2&3/4&9/16&11/32&33/64&35/128
\end{array}
\]
with exponent word $(1,2,1,1,1)$ and color word $(1,1,2,1,2)$.
Its two adjacent internally bracketed runs have colors $2$ and $1$, and
both have carrier $47$:
\[
\begin{aligned}
  2x_2+m_2&=94=2\cdot47,
  &2x_3+m_3&=141=3\cdot47,\\
  x_3+m_3&=94=2\cdot47,
  &x_4+m_4&=141=3\cdot47.
\end{aligned}
\]
Hence any carrier charging argument must use fixed-ray information beyond
finite strict-barrier synchronization.
\end{theorem}

\begin{proof}
For the first segment, direct calculation gives
\[
  2^{a_j}x_{j+1}=3x_j+1,
  \qquad
  2^{a_j}m_{j+1}=3m_j-c_j
\]
with $(c_0,c_1,c_2)=(2,1,3)$.  The displayed rounding errors satisfy
\[
  2^{a_j}\theta_{j+1}=3\theta_j+1-c_j,
\]
and lie in $[0,1)$.  Therefore $e_j:=m_j-\theta_j>1$ and subtraction proves
$2^{a_j}e_{j+1}=3e_j-1$.  The cumulative exponents are $(1,3,4)$, so
\[
  2^1<3,
  \qquad
  2^3<3^2,
  \qquad
  2^4<3^3.
\]
The carrier calculation in the statement proves the first assertion.

For the second segment, the same three direct identities hold with the
displayed exponent, color, and rounding-error words.  Its cumulative
exponents are $(1,3,4,5,6)$, and
\[
  2,8,16,32,64<3,9,27,81,243
\]
coordinatewise.  Thus this word is also a strict barrier.  The two carrier
calculations prove the failure of local injectivity.
\end{proof}

\begin{theorem}[projective defect charge and zero-defect reduction]
Retain the normalized fixed-seed setup, and define the projective rounding
coordinate
\[
  p_j:=\frac{\theta_j}{m_j}.
\tag{114}
\]
On every nonpositive-defect edge, with $c_j=1-d_j\in\{1,2,3\}$, we have
\[
\boxed{
  p_{j+1}-p_j
  =\frac{(1-c_j)m_j+c_j\theta_j}
         {m_j(3m_j-c_j)}.}
\tag{115}
\]
Consequently,
\[
\begin{array}{c|c}
\text{edge type}&\text{projective charge}\\ \hline
c_j=1&
  0\leq p_{j+1}-p_j<\displaystyle\frac1{8m_j},\\[2mm]
c_j=2,\ m_j\geq4&
  \displaystyle\frac1{m_j}<5(p_j-p_{j+1}),\\[2mm]
c_j=3&
  \displaystyle\frac1{m_j}<2(p_j-p_{j+1}).
\end{array}
\tag{116}
\]

Let $\mathcal B$ be the family of maximal consecutive blocks of
nonpositive-defect edges.  For $B=[\alpha,\beta)\in\mathcal B$, put
\[
\begin{aligned}
  \mathcal Z_B&:=\{j\in B:c_j=1\},\\
  \mathcal T_B&:=\{j\in B:c_j=2,\ m_j=2\},\\
  \mathcal R_B&:=B\setminus(\mathcal Z_B\cup\mathcal T_B).
\end{aligned}
\tag{117}
\]
Then
\[
\boxed{
  \sum_{j\in\mathcal R_B}\frac1{m_j}
  \leq\frac5{m_\alpha}
  +\frac58\sum_{j\in\mathcal Z_B}\frac1{m_j}.}
\tag{118}
\]
The same inequality holds for a final infinite block, with the sums
interpreted as monotone limits.

Put
\[
  \mathcal E:=\{\alpha:[\alpha,\beta)\in\mathcal B\},
  \qquad
  \mathcal Z:=\bigcup_{B\in\mathcal B}\mathcal Z_B,
  \qquad
  \mathcal T:=\bigcup_{B\in\mathcal B}\mathcal T_B.
\tag{119}
\]
Since the total nonpositive-defect reciprocal mass diverges, at least one of
the following three series diverges:
\[
\boxed{
  \sum_{j\in\mathcal E}\frac1{m_j},
  \qquad
  \sum_{j\in\mathcal Z}\frac1{m_j},
  \qquad
  \sum_{j\in\mathcal T}\frac1{m_j}.}
\tag{120}
\]

Finally, let $I=[s,t]$ be a color-$1$ run with $\theta_s>0$.  Then
\[
\boxed{
  \ell_I:=\log\frac{p_t}{p_s}
  =\sum_{j=s}^{t-1}-\log\left(1-\frac1{3m_j}\right),
  \qquad
  2\ell_I\leq\sum_{j=s}^{t-1}\frac1{m_j}<3\ell_I.}
\tag{121}
\]
Only finitely many color-$1$ runs on the normalized ray can have
$\theta_s=0$.  Therefore the middle alternative in (120) is equivalent to
\[
  \sum_{I\text{ a color-}1\text{ run}}\ell_I=\infty.
\tag{122}
\]
More generally, define
\[
  P_1(p):=p,
  \qquad
  P_2(p):=1-2p,
  \qquad
  P_3(p):=2-3p.
\tag{123}
\]
These quantities are nonnegative on an edge of the corresponding color.
Whenever $P_c(p_j)>0$, we have
\[
\boxed{
  \frac{P_c(p_{j+1})}{P_c(p_j)}
  =\left(1-\frac{c}{3m_j}\right)^{-1}.}
\tag{124}
\]
Hence, for every color-$c$ run $I=[s,t]$ with $P_c(p_s)>0$,
\[
  L_I:=\log\frac{P_c(p_t)}{P_c(p_s)}
  \qquad\Longrightarrow\qquad
  \frac2cL_I\leq\sum_{j=s}^{t-1}\frac1{m_j}<\frac3cL_I.
\tag{125}
\]
In particular, an $m=2$ tower of $K$ edges satisfies
\[
  \log\frac{P_2(p_{s+K})}{P_2(p_s)}=K\log\frac32,
  \qquad
  \sum_{j=s}^{s+K-1}\frac1{m_j}=\frac K2.
\tag{126}
\]
Thus all divergent nonpositive mass is forced into block entrances,
logarithmic amplification of $P_1$ on zero-defect runs, or logarithmic
amplification of $P_2$ on terminal towers.
\end{theorem}

\begin{proof}
The ceiling and rounding recurrences are
\[
  2^{a_j}m_{j+1}=3m_j-c_j,
  \qquad
  2^{a_j}\theta_{j+1}=3\theta_j+1-c_j.
\]
We divide the second identity by the first and subtract
$p_j=\theta_j/m_j$.  This proves (115).

If $c_j=1$, equation (115) gives
\[
  p_{j+1}-p_j=\frac{\theta_j}{m_j(3m_j-1)}.
\]
Parity gives $m_j\geq3$, and $0\leq\theta_j<1$, which proves the first
line of (116).  If $c_j=2$, then
\[
  p_j-p_{j+1}=\frac{m_j-2\theta_j}{m_j(3m_j-2)}.
\]
For $m_j\geq4$, we have
\[
  \frac{1/m_j}{p_j-p_{j+1}}
  =\frac{3m_j-2}{m_j-2\theta_j}
  <\frac{3m_j-2}{m_j-2}\leq5.
\]
The function in the last member is decreasing for $m_j\geq4$.  Finally,
if $c_j=3$, parity gives $m_j\geq3$, and
\[
  \frac{1/m_j}{p_j-p_{j+1}}
  =\frac{3m_j-3}{2m_j-3\theta_j}
  <\frac{3m_j-3}{2m_j-3}\leq2.
\]
This proves (116).

Fix a block $B=[\alpha,\beta)$.  Let $U_B$ be the total increase of $p_j$
on $\mathcal Z_B$, and let $D_B$ and $D_B^{\mathrm{tower}}$ be its total
decreases on $\mathcal R_B$ and $\mathcal T_B$, respectively.  Telescoping
over a finite initial part of the block gives
\[
  0\leq p_{\mathrm{end}}
  =p_\alpha+U_B-D_B-D_B^{\mathrm{tower}}.
\]
Therefore $D_B\leq p_\alpha+U_B$.  We have
\[
  p_\alpha<\frac1{m_\alpha},
  \qquad
  U_B<\frac18\sum_{j\in\mathcal Z_B}\frac1{m_j}.
\]
The last two lines of (116), with the weaker common constant $5$, now prove
(118).  Passing through finite initial parts proves the infinite-block
case.  Summing (118) over $\mathcal B$ shows that convergence of all three
series in (120) would force convergence on $\bigcup_B\mathcal R_B$ as well.
This would contradict the divergent nonpositive-defect mass.  Hence (120)
holds.

On a color-$1$ edge with $p_j>0$, direct division gives
\[
  \frac{p_{j+1}}{p_j}
  =\frac{3m_j}{3m_j-1}
  =\left(1-\frac1{3m_j}\right)^{-1}.
\]
Multiplication along the run proves the identity in (121).  We also have
\[
  0<\frac1{3m_j}\leq\frac19,
  \qquad
  \frac1{3m_j}
  <-\log\left(1-\frac1{3m_j}\right)
  \leq\frac1{2m_j}.
\]
This proves the two-sided estimate.  If $\theta_s=0$, every color-$1$ edge
preserves $\theta=0$, and the corresponding run anchor equals
$\mathcal K$.  The argument surrounding (79) proves that this equality can
occur only finitely many times on the normalized ray.  Removing those
finitely many runs and summing (121) proves (122).

It remains to prove the unified identity.  For $c=1$, equation (124) is the
ratio already used in (121).  For $c=2$, we have
\[
  2^{a_j}(m_{j+1}-2\theta_{j+1})=3(m_j-2\theta_j),
\]
while for $c=3$ we have
\[
  2^{a_j}(2m_{j+1}-3\theta_{j+1})
  =3(2m_j-3\theta_j).
\]
Division by the corresponding ceiling recurrence proves (124).  We sum
its logarithm along a run and apply the inequality used in (104), which
proves (125).  On an $m=2$ tower, every term in (124) equals $3/2$, and
every reciprocal term equals $1/2$.  This proves (126).
\end{proof}

\begin{theorem}[threshold-reset identity for internal zero-defect runs]
Retain the normalized fixed-seed setup.  Let $I=[s,t]$ be a maximal
color-$1$ run with $\theta_s>0$.  Suppose that it is internally bracketed
by nonpositive-defect edges of colors
\[
  c_-:=c_{s-1}\in\{2,3\},
  \qquad
  c_+:=c_t\in\{2,3\}.
\tag{127}
\]
We define its entrance threshold residual and size-corrected reset depth by
\[
\begin{aligned}
  \rho_I&:=3\theta_{s-1}-(c_--1)
          =2^{a_{s-1}}\theta_s,\\
  \Gamma_I&:=
  \log\frac{3m_{s-1}-c_-}{m_t\rho_I}.
\end{aligned}
\tag{128}
\]
Then $\rho_I>0$ and
\[
\boxed{
  p_s=\frac{\rho_I}{3m_{s-1}-c_-},
  \qquad
  \Gamma_I=\log\frac{m_s}{m_t\theta_s}.}
\tag{129}
\]
The exact projective amplification from (121) satisfies
\[
\boxed{
  \ell_I=\Gamma_I+\log\theta_t,
  \qquad
  0<\Gamma_I-\ell_I
  \leq\log\frac3{c_+-1}.}
\tag{130}
\]
Equality holds in the upper bound precisely when the edge following the
run has $\theta_{t+1}=0$.

Consequently, for every collection $\mathcal I$ of internally bracketed
color-$1$ runs,
\[
  \sum_{I\in\mathcal I}\sum_{j=s(I)}^{t(I)-1}\frac1{m_j}=\infty
  \quad\Longrightarrow\quad
  \sum_{I\in\mathcal I}\ell_I=\infty
  \quad\Longrightarrow\quad
  \sum_{I\in\mathcal I}\Gamma_I=\infty.
\tag{131}
\]
Thus any summable upper charge for the reset depths $\Gamma_I$ closes the
internally bracketed part of the zero-defect branch.
\end{theorem}

\begin{proof}
On the edge entering the run, the ceiling and rounding recurrences give
\[
  2^{a_{s-1}}m_s=3m_{s-1}-c_-,
  \qquad
  2^{a_{s-1}}\theta_s
  =3\theta_{s-1}+1-c_-=\rho_I.
\tag{132}
\]
The assumption $\theta_s>0$ proves $\rho_I>0$.  Dividing the two
identities in (132) proves the first identity in (129), and substituting
them in (128) proves the second.

At the endpoints of the color-$1$ run, we have
\[
\begin{aligned}
  \ell_I
  &=\log\frac{p_t}{p_s}
    =\log\frac{m_s\theta_t}{m_t\theta_s}\\
  &=\Gamma_I+\log\theta_t.
\end{aligned}
\tag{133}
\]
The edge following the run has color $c_+$, so
\[
  2^{a_t}\theta_{t+1}=3\theta_t+1-c_+\geq0.
\tag{134}
\]
Hence
\[
  \frac{c_+-1}{3}\leq\theta_t<1.
\tag{135}
\]
Equations (133)--(135) prove (130), including its equality statement.
Finally, (121) compares the reciprocal mass of each run with $\ell_I$,
while (130) gives $\Gamma_I>\ell_I$.  Summation proves (131).
\end{proof}

\begin{theorem}[strict-barrier failure of local reset bounds]
For every integer $K\geq1$, the exponent word
\[
  w_K:=\bigl(1,(1,2)^K,1\bigr)
\tag{136}
\]
has a positive odd ordinary realization and a synchronized real companion
segment with ceiling and color data
\[
\begin{aligned}
  &(m_0,m_1,\ldots,m_{2K+2})
    =\bigl(4,5,(7,5)^K,6\bigr),\\
  &(c_0,c_1,\ldots,c_{2K+1})
    =\bigl(2,1^{\,2K},3\bigr).
\end{aligned}
\tag{137}
\]
Its middle $2K$ edges form one internally bracketed color-$1$ run
$I_K=[1,2K+1]$.  This run consists of $K$ laps of
\[
  5\xrightarrow{\,a=1\,}7
   \xrightarrow{\,a=2\,}5
\tag{138}
\]
and satisfies
\[
\boxed{
\begin{aligned}
  \sum_{j=1}^{2K}\frac1{m_j}&=\frac{12K}{35},\\
  \ell_{I_K}&=K\log\frac98,\\
  \rho_{I_K}&=\frac43\left(\frac89\right)^K,\\
  \Gamma_{I_K}&=K\log\frac98+\log\frac32.
\end{aligned}}
\tag{139}
\]
Every prefix of $w_K$ obeys the strict-barrier inequality
\[
  2^{A_n}<3^n,
  \qquad
  A_n:=\sum_{j<n}a_j.
\tag{140}
\]

In particular, the complete discrete boundary tuple
\[
  (c_-,a_{s-1},m_{s-1},m_s,m_t,c_+,a_t,m_{t+1})
  =(2,1,4,5,5,3,1,6)
\tag{141}
\]
is independent of $K$, while both the run mass and its projective
amplification tend to infinity.  Therefore no finite-valued bound depending
only on this boundary tuple, or on any fixed-width window of exponent,
color, and ceiling data at the two boundaries, can control zero-defect
mass or projective amplification on finite strict-barrier synchronized
segments.  A valid bound must detect the threshold residual in (128), or
use coherence along one fixed ordinary ray across different runs.
\end{theorem}

\begin{proof}
We first construct the real companion.  Set
\[
  \theta_1:=\frac23\left(\frac89\right)^K,
  \qquad
  \theta_0:=\frac{1+2\theta_1}{3}.
\tag{142}
\]
Use the ceiling path in (137).  The first edge satisfies
\[
  2\cdot5=3\cdot4-2,
  \qquad
  2\theta_1=3\theta_0-1.
\tag{143}
\]
On each color-$1$ lap in (138), the ceiling identities are
\[
  2\cdot7=3\cdot5-1,
  \qquad
  4\cdot5=3\cdot7-1,
\tag{144}
\]
and the corresponding two rounding updates multiply $\theta$ by
$(3/2)(3/4)=9/8$.  Hence
\[
  \theta_{2K+1}=\left(\frac98\right)^K\theta_1=\frac23.
\tag{145}
\]
The final color-$3$ edge satisfies
\[
  2\cdot6=3\cdot5-3,
  \qquad
  2\theta_{2K+2}=3\theta_{2K+1}-2=0.
\tag{146}
\]
All rounding errors lie in $[0,1)$.  Indeed, before the second half of the
last lap the error is $16/27$, after its first half it is $8/9$, and all
earlier errors are smaller.  Also $1/3<\theta_0<1$.  Thus
$e_j:=m_j-\theta_j>1$, and (143)--(146) prove the synchronized real
recurrence.

We next verify the strict barrier.  Put $\lambda:=\log_2 3$.  After the
initial exponent and $k$ complete pairs, the barrier surplus is
\[
  \lambda(1+2k)-(1+3k)
  =\lambda-1+(2\lambda-3)k>0.
\tag{147}
\]
After the first exponent of the next pair, and also after the final extra
exponent, it is
\[
  \lambda(2+2k)-(2+3k)
  =2(\lambda-1)+(2\lambda-3)k>0.
\tag{148}
\]
Here $\lambda>1$ follows from $3>2$, and $2\lambda>3$ follows from
$3^2>2^3$.  Equations (147)--(148) cover every prefix and prove (140).

For completeness, we prove that every finite positive exponent word
$(a_0,\ldots,a_{N-1})$ has a positive odd ordinary realization.  Put
\[
  A_k:=\sum_{j<k}a_j,
  \qquad
  C_k:=\sum_{i=0}^{k-1}3^{k-1-i}2^{A_i}
  \qquad(0\leq k\leq N).
\tag{149}
\]
Iteration of the ordinary recurrence gives
\[
  2^{A_k}x_k=3^kx_0+C_k
  \qquad(0\leq k\leq N).
\tag{150}
\]
Since $3^N$ is invertible modulo $2^{A_N+1}$, there is a unique residue
class $x_0$ modulo $2^{A_N+1}$ satisfying
\[
  3^Nx_0+C_N\equiv2^{A_N}\pmod{2^{A_N+1}}.
\tag{151}
\]
The integer $C_N$ is odd, so this residue class is odd.  We take its least
positive representative.

We check all intermediate valuations.  For $k<N$, split $C_N$ at $k$ and
reduce (151) modulo $2^{A_k+1}$.  The term with index $i=k$ is the only
tail term not already divisible by $2^{A_k+1}$, and we obtain
\[
  3^{N-1-k}
  \left(3(3^kx_0+C_k)+2^{A_k}\right)
  \equiv0\pmod{2^{A_k+1}}.
\tag{152}
\]
Multiplication by odd units does not change the unique nonzero class
$2^{A_k}$ modulo $2^{A_k+1}$.  Therefore
\[
  3^kx_0+C_k\equiv2^{A_k}\pmod{2^{A_k+1}}.
\tag{153}
\]
Together with (151), this proves that every quotient
$x_k:=(3^kx_0+C_k)/2^{A_k}$ is a positive odd integer.  Consecutive
instances of (150) give
\[
  2^{a_k}x_{k+1}=3x_k+1,
\]
with the prescribed exact exponent.  Applying this construction to $w_K$
proves the ordinary realization asserted in the theorem.

It remains to calculate the charges.  Each lap contributes
$1/5+1/7=12/35$ to the reciprocal mass.  Its two rounding updates multiply
$p$ by $9/8$, because the initial and terminal ceiling of a lap are both
$5$.  This proves the first two identities in (139).  Equations
(142)--(143) give
\[
  \rho_{I_K}=3\theta_0-1=2\theta_1
  =\frac43\left(\frac89\right)^K.
\tag{154}
\]
Finally, $m_s=m_t=5$ and (129) give
\[
  \Gamma_{I_K}
  =\log\frac1{\theta_1}
  =K\log\frac98+\log\frac32.
\tag{155}
\]
This proves (139).  The tuple (141) is read directly from (136)--(137).
For any fixed boundary-window width, taking $K$ larger than that width
leaves the two discrete boundary windows unchanged while (139) grows
without bound.  The claimed impossibility follows.
\end{proof}

\begin{theorem}[global logarithmic separation clock]
Retain the normalized fixed-seed setup and define
\[
  \mathcal L_j:=\log\frac{q_j}{m_j}.
\tag{156}
\]
For every edge, with $c_j:=1-d_j$, one has
\[
\boxed{
  \mathcal L_{j+1}-\mathcal L_j
  =\log\frac{3m_j}{3m_j-c_j}
  =-\log\left(1+\frac{d_j-1}{3m_j}\right).}
\tag{157}
\]
The sign of this increment is determined exactly by the defect:
\[
\begin{array}{c|c|c}
d_j&c_j&\mathcal L_{j+1}-\mathcal L_j\\ \hline
d_j\leq0&c_j\in\{1,2,3\}&>0\\
d_j=1&c_j=0&=0\\
d_j\geq2&c_j\leq-1&<0.
\end{array}
\tag{158}
\]
On a nonpositive-defect edge,
\[
\boxed{
  \frac2{c_j}(\mathcal L_{j+1}-\mathcal L_j)
  \leq\frac1{m_j}
  <\frac3{c_j}(\mathcal L_{j+1}-\mathcal L_j).}
\tag{159}
\]
Consequently, every consecutive nonpositive-defect block
$B=[\alpha,\beta)$ satisfies
\[
\boxed{
  \frac23(\mathcal L_\beta-\mathcal L_\alpha)
  \leq\sum_{j=\alpha}^{\beta-1}\frac1{m_j}
  <3(\mathcal L_\beta-\mathcal L_\alpha).}
\tag{160}
\]
The same assertion holds for a final infinite block, using monotone
limits.

Let $I=[s,t]$ be a constant nonpositive-defect run of color
$c\in\{1,2,3\}$, and let $\Lambda_I$ be its rational anchor from (70).
Put
\[
  \Delta_I:=
  \begin{cases}
    \Lambda_I-\mathcal K,&c=1,\\
    c\mathcal K-\Lambda_I,&c\in\{2,3\}.
  \end{cases}
\tag{161}
\]
Whenever $\Delta_I>0$, the projective potential in (123) has the exact
factorization
\[
\boxed{
  P_c(p_j)=\Delta_I\frac{q_j}{m_j}
  =\Delta_I e^{\mathcal L_j}
  \qquad(s\leq j\leq t).}
\tag{162}
\]
In particular, the projective run gain in (125) is
\[
\boxed{L_I=\mathcal L_t-\mathcal L_s.}
\tag{163}
\]
For an internally bracketed color-$1$ run with positive entrance rounding,
the quantities in the threshold-reset theorem therefore satisfy
\[
\boxed{
  \ell_I=\mathcal L_t-\mathcal L_s,
  \qquad
  \Gamma_I=\mathcal L_t-\mathcal L_s-\log\theta_t.}
\tag{164}
\]

Define the accumulated nonpositive gain and positive-defect contraction by
\[
\begin{aligned}
  \mathcal C_N
  &:=\sum_{\substack{0\leq j<N\\d_j\leq0}}
    -\log\left(1+\frac{d_j-1}{3m_j}\right),\\
  \mathcal P_N
  &:=\sum_{\substack{0\leq j<N\\d_j\geq2}}
    \log\left(1+\frac{d_j-1}{3m_j}\right).
\end{aligned}
\tag{165}
\]
Then
\[
\boxed{
  \mathcal L_N-\mathcal L_0
  =\mathcal C_N-\mathcal P_N,}
\tag{166}
\]
and
\[
  \frac23\mathcal C_N
  \leq\sum_{\substack{0\leq j<N\\d_j\leq0}}\frac1{m_j}
  \leq3\mathcal C_N.
\tag{167}
\]
Moreover,
\[
\boxed{
\begin{aligned}
  &\log\frac{x_{j+1}/m_{j+1}}{x_j/m_j}\\
  &\qquad=\mathcal L_{j+1}-\mathcal L_j
  +\log\left(1+\frac1{3x_j}\right).
\end{aligned}}
\tag{168}
\]
Since $\sum_j1/x_j<\infty$, the last terms in (168) have a finite total.
Thus
\[
\boxed{
  \sum_{d_j\leq0}\frac1{m_j}=\infty
  \quad\Longleftrightarrow\quad
  \sum_{d_j\leq0}
  \log\frac{x_{j+1}/m_{j+1}}{x_j/m_j}=\infty.}
\tag{169}
\]
Finally, $m_j/q_j\to0$, so $\mathcal L_j\to\infty$.  Hence every
normalized divergent ray obeys the strict drift law
\[
\boxed{\mathcal C_N-\mathcal P_N\longrightarrow\infty.}
\tag{170}
\]
Therefore all divergent nonpositive reciprocal mass, including zero-defect
cycle laps and terminal $m=2$ towers, is quantitatively equivalent to the
positive variation of the single clock $\mathcal L_j$.  A fixed-ray closure
must contradict (170), not merely bound one color in isolation.
\end{theorem}

\begin{proof}
The ceiling recurrence is
\[
  2^{a_j}m_{j+1}=3m_j-c_j,
\]
while $q_{j+1}/q_j=3/2^{a_j}$.  Division proves (157), and the three signs
in (158) follow immediately.

Suppose that $d_j\leq0$.  Then $c_j\in\{1,2,3\}$.  Put
\[
  z:=\frac{c_j}{3m_j}.
\]
Parity in the ceiling recurrence gives $m_j\geq c_j$, and hence
$0<z\leq1/3$.  The elementary estimates
\[
  z<-\log(1-z)\leq\frac{z}{1-z}\leq\frac32z
\]
now prove (159).  Summing over a block and using
$1\leq c_j\leq3$ proves (160).

Equation (72) gives
\[
  \Lambda_I-c\mathcal K
  =\frac{c\theta_j-(c-1)m_j}{q_j}.
\]
For $c=1$, this says
$P_1(p_j)=\theta_j/m_j=(\Lambda_I-\mathcal K)q_j/m_j$.
For $c\in\{2,3\}$, it says
\[
  P_c(p_j)=(c-1)-cp_j
  =(c\mathcal K-\Lambda_I)\frac{q_j}{m_j}.
\]
This proves (162).  Taking endpoint ratios proves (163), and (164) follows
from (130).

Equations (166) and (167) are obtained by summing (157) and (159), noting
that a defect $d_j=1$ contributes zero.  Next, $h_j=x_j/q_j$ and (98) give
\[
  \frac{h_{j+1}}{h_j}=1+\frac1{3x_j}.
\]
Since
\[
  \log\frac{x_j}{m_j}=\mathcal L_j+\log h_j,
\]
equation (168) follows.  The inequality
$\log(1+u)\leq u$ and the convergence of $\sum_j1/x_j$ show that the total
correction in (168) is finite.  On every nonpositive edge, both terms on
the right of (168) are positive.  Combining this observation with (167)
proves (169).

Finally, equation (78) gives $m_j/q_j\to0$.  Therefore
$\mathcal L_j=\log(q_j/m_j)\to\infty$, and (166) proves (170).
\end{proof}

\begin{corollary}[upper-anchor defect-tail identity]
For every state index $j$, the ordered defect tail converges and satisfies
\[
\boxed{
  \frac{\theta_j}{q_j}
  =\Omega_j-\mathcal K
  =-\sum_{k=j}^{\infty}\frac{d_k}{3q_k}.}
\tag{171}
\]
Let $I=[s,t]$ be an internally bracketed color-$1$ run as in the
threshold-reset theorem.  Then
\[
\boxed{
  \rho_I=q_{s-1}
  \left(
    \frac{c_+-1}{q_t}
    -\sum_{k=t+1}^{\infty}\frac{d_k}{q_k}
  \right).}
\tag{172}
\]
Equivalently,
\[
\boxed{
  \rho_I=q_{s-1}
  \left(
    \frac{c_+-1}{q_t}
    +\frac{3\theta_{t+1}}{q_{t+1}}
  \right).}
\tag{173}
\]
Thus the entrance residual is exactly the exit threshold jump, transported
back to the entering scale, together with the complete future defect tail.
\end{corollary}

\begin{proof}
Equation (53) gives
\[
  \Omega_{k+1}-\Omega_k=\frac{d_k}{3q_k}.
\]
Also $\Omega_k-\mathcal K=\theta_k/q_k\to0$.  Telescoping from $j$ to
$N$ and letting $N\to\infty$ proves (171), without any assertion of
absolute convergence.

On the run, $d_k=0$ for $s\leq k<t$, while
$d_t=1-c_+=-(c_+-1)$.  Substitution in (171), followed by
$\rho_I=3q_{s-1}(\Omega_s-\mathcal K)$ from (128)--(129), proves (172).
Applying (171) at $t+1$ gives
\[
  \sum_{k=t+1}^{\infty}\frac{d_k}{q_k}
  =-\frac{3\theta_{t+1}}{q_{t+1}},
\]
which proves (173).
\end{proof}

\end{document}
